\section{Non-stationary Reinforcement Learning}

In this section, we formalise the notion of non-stationarity. 
We then present related approaches to dealing with it, particularly in the context of lifelong learning.   


\subsection{The Non-stationarity in Reinforcement Learning}
\label{subsec:nature_ns}
\subsubsection{Nature of Non-stationarity}
In \gls{rl}, 4 elements of the \gls{mdp} out of the five defining the process can induce non-stationarity. 
Obviously, the reward and transition functions might depend on time. 
More occasionally, the state space or the action space itself can vary over time. 
A further difference can be made between processes whose non-stationarity is present in between episodes and those where it can happen inside an episode. 
We refer to the first as \emph{inter-episode} non-stationarity. 
It can be thought of as a set of tasks from which a task is selected at each episode. 
We refer to the second type as \emph{intra-episode} non-stationarity. 
In this case, the notion of task is less clear as the dynamics of the \gls{mdp}  can change at each decision step. 
Therefore, the estimation and analysis of past data must be even more careful.
Lastly, it is usually considered that the non-stationarity can be either \emph{smooth} when the dynamics evolve with some regularity through time or \emph{abrupt} when the dynamics typically change less often but with a clearer shift in behaviour. 

Concerning the peculiar case of delay-induced non-stationarity, due to the previously mentioned partial observability induced by the memoryless approach, the transition dynamics and the reward will be the
source of non-stationarity. 
Of course, the state and action spaces will remain stationary. 
The non-stationarity is by nature intra-episode, as the dynamics will depend on the unobserved buffer of actions inside the augmented state.
Finally, the non-stationarity will be smooth, as the augmented state's action buffer shares many actions with the following augmented state.
Note that, as the process is divided into discrete time steps, even smooth non-stationarity can be seen as abrupt in between time steps. 
We will explain more formally how we define the smoothness of a discrete process in \Cref{subsec:bias_analysis}.


\subsubsection{Drivers of Non-stationarity}
The non-stationarity mentioned in the previous section might have different causes. We present three of them, following the taxonomy of \cite{khetarpal2020towards}.\\
\noindent\textbf{Multi-task Setting.}\indent Arguably the most common setting in the non-stationary \gls{rl} literature, this setting considers a set $(\markovdecisionprocess_i)_{i\in\naturalnumbers}$ of stationary \glspl{mdp}--or tasks--where the environment might regularly switch from one task to another during the learning and testing phases.
\\ 
\noindent\textbf{Passive Non-stationarity.}\indent This setting considers a type non-stationarity that occurs irrespective of the agent's actions.
\\
\noindent\textbf{Active Non-stationarity.}\indent Unlike in the previous setting, here, the agent may affect the non-stationarity of the environment.

The delay induces an active non-stationarity since it is the agent that chooses the actions in the replay buffer, and it is the partial knowledge induced by ignoring these actions that cause the non-stationarity.

\subsubsection{Formalisation of Lifelong Reinforcement Learning}
In this sub-section, we formalise the framework of lifelong \gls{rl} that we will be interested in. 
It is based on a non-stationary \gls{mdp} $\markovdecisionprocess$ with transition functions $(\transitionfunction_t)_{t \in \naturalnumbers}$ and reward functions 
$(\rewardfunction_t)_{t \in \naturalnumbers}$ of expected values $(\expectedrewardfunction_t)_{t \in \naturalnumbers}$.
These functions define the transition and the reward at each decision step $t \in \naturalnumbers$. 
We assume that $\lVert r_t\rVert_{\infty} \le \maximumreward < \infty $ for any $t\in\naturalnumbers$ and consider a discount factor $\gamma$ in $[0,1]$.
We then define a non-stationary policy $\pi = (\pi_t)_{t \in \naturalnumbers}$ that defines the policy that the agent follows at each decision step $t\in\naturalnumbers$. 

The goal in lifelong \gls{rl} differs from the one in traditional \gls{rl} since the underlying process cannot be reinitialised.
The environment is initialised once and for all, and the agent cannot collect more than a single trajectory of experience. 
For example, consider the problem of trading, where it is clear that the exact same market conditions will never be met twice. 
Without an accurate market simulator, the agent cannot go back in time and try another strategy. 
On the contrary, an example of a traditional \gls{rl} setting is chess. 
The agent can play several games, always starting from the same initial state distribution, and it can therefore gather more information about similar environment conditions. 
In lifelong \gls{rl}, the expected future return, therefore, depends on the agent's current state. 
This leads to the following objective called \emph{$\beta$-steps ahead expected return}.
Let $T \in \naturalnumbers$ be the current time and let $\beta \in \naturalnumbers$, the {$\beta$-steps ahead expected return} is defined as:
\begin{align}\label{eq:beta_steps_ahead_return}
	\expectedreturn[T,\beta][](\pi) =\sum_{t = T+1}^{T+\beta} \widehat{\gamma}^t \expectedvalue_t^\pi [\expectedrewardfunction_{t}(s_t,a_t)],
\end{align}
where $\widehat{\gamma}^t = \gamma^{t-T-1}$ and $\mathbb{E}_t^\pi$ is the expectation under the state-action distribution induced by the non-stationary policy $\pi$ after $t$ steps.
For ease of notation, in the following, we remove the dependence on the pair $(s_t,a_t)$ of the reward from the notation. 

Finally, as for classic \glspl{mdp}, we say that a policy $\optimalpolicy$ is $\beta$-step ahead optimal at time $T$ if 
\begin{align*}
    \optimalpolicy\in \argmax_{\pi \in \Pi^{\text{HR}}} \expectedreturn[T,\beta][](\pi),
\end{align*}
where $\Pi^{\text{HR}}$ is the set of non-stationary policies as defined in \Cref{subsec:policies}.

We could draw a parallel and say that the goal of classical \gls{rl} is to maximise $\expectedreturn[0,H][]$ where $H$ is the (possibly infinite) horizon. 
However, in this case, the agent might restart the process multiple times to collect different episodes of this same process.
Instead, a lifelong learning agent is interested in maximising 
$\expectedreturn[T,\infty][]$ after having experienced exactly $T$ transitions of a single episode. 


\subsection{Related Works in Reinforcement Learning for Non-stationarity}
\label{subsec:related_worls_lifelong}

Adapting \gls{rl} to non-stationary \glspl{mdp} has been widely studied in the literature~\cite{garcia2000solving, ghate2013linear, lesner2015nonstationary}. 
In order to adapt to new tasks, the agent must learn about the structure of the overall problem. 
For instance, the agent can decompose the problem into smaller sub-problems using function composition \cite{griffiths2019subtasks} or learn a low-dimensional abstract representation of the environments \cite{zhang2018decoupling, Francois-Lavet2019abstract}.
Another related approach is to focus not on task-specific dynamics but rather on task-agnostic underlying dynamics of the problem in order to learn an efficient general policy. 
This can be achieved by building auxiliary tasks such as reward prediction \cite{jaderberg2017rew-pred} or by inverse dynamics prediction \cite{Shelhamer2017dyn-pred}.

Although stationary by nature, meta \gls{rl} can be adapted to non-stationary processes. 
Solutions by meta \gls{rl} can be interesting, as they use previous experience to learn new skills efficiently.
Essentially, adapting these algorithms to non-stationary processes consists of casting back the problem to a stationary one.  
For instance, the sequence of tasks can be modelled as a Markov chain \cite{al-shedivat2018continuous}, an experience replay buffer of tasks can be used to make the process ``more stationary'' \cite{riemer2018learning} as is done usually in \gls{rl} \cite{mnih2013playing} or a latent model of the succession of tasks can be learnt \cite{poiani2021meta}.

The previous works usually apply to inter-episode non-stationarity. 
We will now study the problem of intra-episode non-stationarity, where dynamics can change inside an episode, such as in lifelong \gls{rl}. 
In finite-horizon settings, with some modifications, theoretical \gls{rl} solutions can provide guarantees even in the face of a changing environment \cite{hallak2015contextual,ortner2020uclr2}.
In the infinite-horizon setting, \ie lifelong learning, different statistical tools borrowed from the stochastic processes and time series literature can be used.
Change detection can be used under the assumption of abrupt changes in dynamics \cite{daSilva2006context, hadoux2014sequential}. 
Hidden-state models can be used to describe the sequence of tasks when it is assumed that the environment might switch between a known number of stationary dynamics \cite{choi2000hidden}. 
If the number of stationary dynamics is unknown, a solution is to learn a factored representation of the policy \cite{mendez2020lifelong}. 
In the proposed solution, the first factor of the policy is  
shared on all tasks and is therefore trained to perform well on the distribution of tasks observed thus far. The second factor is task-specific and adapts the policy to the current task. 
Thanks to this structure, the algorithm, LPG-FTW, can adapt quickly to new tasks while avoiding catastrophic forgetting.
More in-depth reviews can be found in  \cite{padakandla2021survey} for non-stationary in \gls{rl} and in \cite{khetarpal2020towards} for lifelong/continual \gls{rl}.

We finish this overview by studying more deeply an approach that is more similar to our POLIS algorithm, highlighting the similarities and disparities between them. 
\cite{chandak2020optimizing} take the rather direct but efficient approach of directly optimising for the future predicted performance. 
For a given set of policy parameters $\boldsymbol\theta\in\Theta$, the past performance of a policy $\pi_{\boldsymbol\theta}$ is estimated by importance sampling.
These past performances are then used inside a regression algorithm, which can be leveraged to predict future performances.
If all these steps are differentiable, then optimising the future predicted performance over $\Theta$ by gradient ascent is possible. 
The authors propose two practical algorithms, Pro-OLS and Pro-WLS, which forecast performance using ordinary least-squares regression and weighted least-squares, respectively. 
The latter builds on the link between weighted least-squares regression and weighted importance sampling \cite{hachiya2012importance,mahmood2014weighted} in order to reduce the variance of the estimates due to importance sampling at the expense of adding some bias.
We will now delve into the differences with respect to our approach.
\\
\noindent\textbf{Setting.}\indent A major difference is that \cite{chandak2020optimizing} designed a solution for episodic RL with inter-episode non-stationarity. We consider instead a truly lifelong framework with intra-episode non-stationarity.
\\
\noindent\textbf{Objective.}\indent We also use importance sampling to estimate future performance, although our objectives differ greatly. 
We add a new discounting parameter to control the bias due to non-stationarity, and we add another two terms to the objective, an estimator of the past performance and a variance regularisation term.
\\
\noindent\textbf{Anytime.}\indent Our approach can be retrained at any time during the lifelong episode. Should a significant shift in dynamics be detected, our policy optimisation can be applied immediately in order to adapt to these new dynamics. 
The algorithm can also be updated less often when the non-stationarity is highly foreseeable from the past. This makes our algorithm an anytime algorithm. In contrast, \cite{chandak2020optimizing}'s approach can train the policy in between episodes but keeps it constant inside an episode.
\\
\noindent\textbf{Parameter-based.}\indent We consider the parameter-based policy optimisation setting compared to \cite{chandak2020optimizing}'s action-based policy optimisation setting.\\
